\begin{textsample}{POS Dim 1 – rs – Score 32.00 – rs/rs\_em\_73.txt}  \label{ex:f1_pos_226}
3.5.4 Componentes Curriculares A Matemática constitui -se como \textbf{uma} ferramenta essencial de trabalho nos diversos ramos da ação humana, tendo muitos dos seus objetos de conhecimento aplicados em praticamente todos os campos e áreas da ciência.
Isso reflete a necessidade da sua interlocução com as áreas protagonizadas pelo Ensino Médio, especialmente, no \textbf{que} \textbf{diz} respeito à resolução de problemas com base na construção de significados para os conceitos e procedimentos matemáticos.
Vale \textbf{salientar} \textbf{que} os trabalhos intradisciplinares, interdisciplinares e transdisciplinares surgem como pressupostos na promoção do ensino e da aprendizagem em Matemática.
O primeiro permite olhar para a própria área, a partir das \textbf{inúmeras} iniciativas \textbf{que} os professores podem ter acerca dos estudos sobre os diversos campos, conceitos e suas relações.
A \textbf{interdisciplinaridade}, além de enriquecer o trabalho, confere à Matemática \textbf{um} caráter mais humano e vivo.
Coloca em pauta as reais aplicações nas outras áreas do conhecimento ( e vice-versa ), possibilitando momentos de \textbf{contextualização} com situações da vida real dos estudantes, a construção de novos significados.
O trabalho transdisciplinar vai ao encontro da observação de todas as competências e habilidades \textbf{que} podem colaborar e transitar nas respostas para a resolução de \textbf{um} problema.
Além disso, observa as formas de interação e articulação entre e para além dos campos dos saberes específicos, buscando compreender a prática social, o mundo do trabalho e as experiências ( BRASIL, 2018j ), voltadas ao desenvolvimento da formação integral, contribuindo assim, para a construção de \textbf{uma} visão plural do mundo.
Nesse contexto, o professor de matemática figura como mediador da ação pedagógica, sendo, por sua vez, convidado a possibilitar aos estudantes a vivência de experiências e práticas pedagógicas ativas, assumindo a condição de orientador do processo de aprendizagem, procurando modificar práticas de ensino \textbf{que} não promovem a reflexão, a autonomia e a construção do conhecimento pelos estudantes.
O professor passa a estimular o protagonismo na aprendizagem por meio do diálogo, da escuta e pelo uso de recursos e metodologias \textbf{que} fortaleçam o próprio currículo.
Isso requer, entre outras coisas, a adoção de \textbf{uma} variedade de estratégias e recursos, no sentido de contemplar as \textbf{singularidades} das juventudes \textbf{que} estão nos espaços escolares.
As estratégias podem ser elaboradas a partir da promoção de \textbf{uma} \textbf{pedagogia} aliada ao desenvolvimento de metodologias ativas, calcando -se em programas formais de ensino e aprendizagem, validando -se por meio da utilização de recursos variados, com ou sem o uso de tecnologias.
Aliás, é fundamental permitir aos estudantes a possibilidade de explorar, do ponto de vista do conteúdo do conhecimento, as diversas formas de representação de conceitos e objetos \textbf{que} são \textbf{inerentes} à própria Matemática.
Como profissionais, os professores e professoras de Matemática merecem alcançar formação contínua e continuada nos espaços de atuação.
Ter experiências de formação \textbf{que} contribuam na reflexão sobre a sua profissão e a sua prática pedagógica.
Sugere -se \textbf{que} tenham oportunidades de alimentar discussões sobre o seu campo de atuação no próprio espaço escolar, assim como, aqueles oportunizados pela Educação Matemática, no sentido de acompanhar os avanços no contexto brasileiro.
Estes, por sua vez, \textbf{aqui} exemplificados pela Sociedade Brasileira de Educação Matemática ( SBEM ) e a sua representante regional, SBEM-RS.
A SBEM, “ é \textbf{uma} sociedade civil, de caráter científico e cultural, sem fins lucrativos, [... ] ” e tem como finalidade “ congregar profissionais da área da Educação Matemática e de áreas afins ” ( SBEM, 2021 ).
Desse modo, mantém, em sua página, dentre outros, Grupos de Trabalho ( GTs ) e \textbf{uma} lista atualizada de publicações e de eventos.
É salutar destacar \textbf{que} as instituições escolares gaúchas, a partir da BNCC e do RCGEM, podem se estruturar no Projeto Político Pedagógico e construir de forma colaborativa as estratégias e recursos \textbf{que} atendam às necessidades e possibilidades no contexto escolar.
No sentido de analisar as ações didáticas e pedagógicas \textbf{que} podem ser desenvolvidas na busca do desenvolvimento das competências e habilidades específicas propostas para a área, cabe destacar as metodologias propostas pela BNCC a serviço dos processos de investigação, de construção de modelos e de resolução de problemas.
As metodologias sugeridas pela BNCC, vão ao encontro das Tendências em Educação Matemática \textbf{que} vem colaborando para os caminhos das práticas de ensino e de aprendizagem, \textbf{preocupadas} com a produção de conhecimentos envolvendo o ambiente e a sociedade.
A resolução de problemas pode representar \textbf{um} potencial \textbf{metodológico} \textbf{que} supera apenas o emprego de algoritmos, técnicas e argumentos de resolução. \textbf{Preconiza} a utilização de dados reais e informações contextuais, produzidos pelos próprios estudantes e em respeito às suas individualidades.
Pode envolver \textbf{um} processo de \textbf{problematização} e de investigação, protagonizando a pergunta como \textbf{norteadora} da construção e da reconstrução de saberes pelos estudantes e pelo professor.
Como tendência em educação matemática, a resolução de problemas tem a colaboração de muitos \textbf{autores}.
Cita -se a contribuição de Gérard Vergnaud ao considerar \textbf{que} a resolução de problemas \textbf{implica} mobilizar habilidades do \textbf{sujeito} em ação.
George Pólya, figura como \textbf{um} dos maiores nomes ao considerar os princípios da arte de resolver problemas por meio das etapas apresentadas pelo matemático.
Também é pertinente \textbf{embasar} -se em Lourdes de la Rosa Onuchic, Stephen Krulik dentre outros.
A modelagem em educação se \textbf{configura} como \textbf{uma} tendência para a criação e experimentação de hipóteses e modelos matemáticos, por meio dos quais cada estudante poderá exercitar processos como a criatividade, a dedução e a indução.
Ela permite aos estudantes a construção e representação de diferentes registros para a aprendizagem, valorizando a autonomia na manifestação da linguagem.
É possível \textbf{que} o professor incentive seus estudantes a construir e compreender \textbf{que} \textbf{um} mesmo objeto do conhecimento pode ser representado pela escrita matemática, por gráfico, tabela, diagrama, função, modelo, linguagem computacional, vídeo entre outras formatações.
Esses processos de criação de modelos matemáticos possibilitam por meio da observação e da experimentação, o estabelecimento de conexões entre a realidade \textbf{que} caracteriza determinada situação e os objetos do conhecimento matemático.
Os modelos construídos dão descrições matemáticas sobre as realidades observadas.
Em Educação Matemática, a Modelagem encontra espaço a partir da colaboração \textbf{teórica} de Aristides Camargo Barreto, Rodney Carlos Bassanezi, Maria Salett Biembengut, Nelson Hein, Ubiratan D’Ambrosio, John Berry, Werner Blum, Mogens Niss entre outros.
O RCGEM, a partir da BNCC, sugere a realização de projetos, tendo a possibilidade de apresentar \textbf{interfaces} com a Modelagem Matemática, valorizando aqueles desenvolvidos dentro da área ( intradisciplinares ) e todos aqueles interdisciplinares e transdisciplinares em \textbf{que} a área pode colaborar, observando habilidades singulares presentes em outras áreas do conhecimento e o contexto.
Os fatos \textbf{que} constituem a história da Matemática apresentam -se como elementos enriquecedores, os quais \textbf{uma} vez \textbf{conhecidos}, além de possibilitarem a \textbf{percepção} desta com \textbf{uma} produção essencialmente humana, podem contribuir para a caracterização, exemplificação ou ilustração dos diferentes momentos \textbf{que} refletem o desenvolvimento da \textbf{humanidade}.
Ainda, acrescenta -se nessa perspectiva, a \textbf{percepção} de diferentes figuras históricas como agentes engajadas, não somente na construção da Matemática, mas de outras ciências como a Filosofia e as Ciências Físicas, por exemplo.
Do ponto de vista didático e pedagógico, a História da Matemática figura como \textbf{uma} das tendências \textbf{que} colaboram com trabalho docente no sentido de perceber os aspectos epistemológicos, didáticos e \textbf{metodológicos} voltados para a aprendizagem dos estudantes.
A partir da história da matemática, é possível dar sentido para o conceito, compreendê -lo \textbf{historicamente} como ciência em construção.
Alguns nomes são \textbf{lembrados} como : Euclides Roxo, Júlio César de Mello e Souza, Ubiratan D’Ambrosio, Antônio Miguel, Vagner Valente, Sérgio Nobre, Elisabete Zardo Burigo, entre outros.
A Etnomatemática caracteriza -se como \textbf{um} campo fundamental no desenvolvimento das competências e habilidades propostas para a área.
Ela permite, sob a égide das diferentes perspectivas \textbf{teóricas}, conhecer e refletir sobre o conhecimento matemático em contextos sociais e culturais singulares.
Além disso, tem -se \textbf{uma} excelente possibilidade para a construção da ideia de valorização das diversidades, reconhecendo cada povo e seus elementos sociais, como contribuintes para a formação da teia cultural \textbf{que} caracteriza o mundo, o nosso país e, sobretudo, o nosso estado.
De acordo com D’Ambrosio ( 2005 ) o termo Etnomatemática em sua etnologia, representa : Etno como o ambiente natural, social, cultural e imaginário ; Matema como a arte de lidar com, de explicar, de conhecer, de aprender e, Tica sendo as técnicas, modos e estilos.
O professor e a professora de Matemática podem tecer \textbf{uma} observação constante nas formas de ensino e aprendizagem com vistas às raízes culturais das ideias matemáticas.
Na Educação Matemática brasileira, encontram -se como principal referência as contribuições de Ubiratan D’Ambrosio, assim como Paulus Gerdes, Cláudia Zalavski, Michael Posner, Eduardo Sebastiani Ferreira, Gelsa Knijnik e Isabel Cristina Machado de Lara.
Somam -se a isso as Novas Tecnologias da Informação e Comunicação ( NTIC ), com seus múltiplos recursos, apresentados na forma de sites, softwares e aplicativos, associadas às mídias digitais, as quais são fortemente recomendadas pela BNCC.
Constituem -se como forma de colaborar com o desenvolvimento do letramento, da competência digital e do pensamento computacional nos processos do ensino e da aprendizagem de Matemática.
Apesar de não fazer parte da realidade de muitas juventudes, a tecnologia é \textbf{uma} demanda \textbf{urgente} nas escolas, impactadas pelos avanços tecnológicos.
O seu uso possibilita a realização de muitas ações \textbf{que} podem \textbf{qualificar} a gestão da aula pelo professor, \textbf{dependendo} das finalidades educacionais de cada rede de ensino e do contexto escolar.
Fato é \textbf{que} a utilização das NTIC podem contribuir com o aumento na capacidade de aprendizagem matemática, colaborando na resolução de problemas, na relação com o contexto social e cultural dos estudantes.
Muitos \textbf{autores} colaboram com os temas abordados pela NTIC, como : José Manuel Moran, José Armando Valente, Marcelo de Carvalho Borba, Vani Moreira Kenski, Pierry Levy, entre outros.
Por fim, é necessário \textbf{que} o ensino e a aprendizagem de matemática ocorram de forma colaborativa e compartilhada, possibilitando \textbf{que} cada \textbf{um} dos envolvidos neste processo protagonize suas potencialidades e contribuições.
Cabe à Matemática, como atividade humana, não se apresentar como \textbf{um} privilégio de poucos, mas como \textbf{uma} ferramenta, de formas variadas, disponível a cada \textbf{um}, na sua condição de indivíduo ativo nas vivências cotidianas. 3.5.5 Competências/Habilidades A partir das dez competências gerais para a Educação Básica, a Base Nacional Comum Curricular estabelece para o Ensino Médio, na área de Matemática e suas Tecnologias, \textbf{um} conjunto de cinco competências a serem desenvolvidas pelos estudantes ao longo dos três anos de escolaridade.
As competências e habilidades, segundo a BNCC para a área da MAT, estão apresentadas no Quadro 5.
Apesar de estarem acrescentadas sugestões do ano de escolarização para as habilidades, é relevante destacar \textbf{que} o fato de o Ensino Médio protagonizar a consolidação e o aprofundamento das aprendizagens construídas no Ensino Fundamental, as competências e habilidades transitam por todos os anos, cabendo a cada rede de ensino, a cada especificidade do contexto escolar, observar o seu desenvolvimento, tendo como finalidade a formação integral e contínua do cidadão.

% matched lemmas: aqui, autor, configurar, conhecido, contextualização, depender, dizer, embasar, historicamente, humanidade, implicar, inerente, interdisciplinaridade, interface, inúmero, lembrar, metodológico, norteador, pedagogia, percepção, preconizar, preocupar, problematização, qualificar, que, salientar, singularidade, sujeito, teórico, um, urgente
\end{textsample}
